\documentclass[12pt,a4paper]{article}

\usepackage[czech]{babel}
\usepackage[T1]{fontenc}
\usepackage[utf8]{inputenc}

\usepackage{geometry}
\usepackage[most]{tcolorbox}

\geometry{margin=2.5cm}

\title{\textbf{Rozbor díla „Havran“}}
\author{Pavlo }

\begin{document}
\maketitle

%------------------------------------------------
\section*{Autor}

\begin{tcolorbox}[breakable]

\textbf{Autor:}
\begin{itemize}
\item Edgar Allan Poe
\end{itemize}

\textbf{Název díla:}
\begin{itemize}
\item Havran
\end{itemize}

\textbf{Literární směr:}
\begin{itemize}
\item Romantismus
\end{itemize}

\textbf{Století tvorby:}
\begin{itemize}
\item 1. pol. 19. st. (1850)
\end{itemize}

\textbf{Další autoři stejného směru + dílo:}
\begin{itemize}
\item George Gordon Byron – Childe Haroldova pouť
\item Victor Hugo – Legenda věků
\item Alexander Sergejevič Puškin – Evžen Oněgin
\end{itemize}

\textbf{Další autorova díla:}
\begin{itemize}
\item Jáma a kyvadlo
\item Maska červené smrti
\end{itemize}

\textbf{Specifický styl autora:}
\begin{itemize}
\item Jeho díla jsou silně ovlivněna jeho znalostmi kosmogonie a přírodních věd. Zároveň byl velmi inteligentní.
\end{itemize}

\end{tcolorbox}

%------------------------------------------------
\section*{Charakteristika uměleckého textu}

\begin{tcolorbox}[breakable]

\textbf{Literární druh, žánr, forma:}
\begin{itemize}
\item Poezie, lyrika, balada
\end{itemize}

\textbf{Typ vypravěče:}
\begin{itemize}
\item Ich-forma, lyrický subjekt
\end{itemize}

\textbf{Dominantní slohový postup:}
\begin{itemize}
\item Úvahový a popisný
\end{itemize}

\textbf{Vysvětlení názvu díla:}
\begin{itemize}
\item Název odkazuje na havrana, který v básni symbolizuje smutek, smrt a nezvratnost ztráty.
\end{itemize}

\textbf{Aktuálnost díla:}
\begin{itemize}
\item Dílo je stále aktuální, protože zpracovává témata smutku, samoty, ztráty a lidské psychiky.
\end{itemize}

\textbf{Místo a čas děje:}
\begin{itemize}
\item V pokoji básníka, pravděpodobně v jeho domě; děj se odehrává v noci.
\end{itemize}

\textbf{Stručné nastínění děje:}
\begin{itemize}
\item Lyrický subjekt sedí doma s horečkou a čte staré svazky, když náhle uslyší zaklepání na dveře. Napadne ho jeho zesnulá milovaná Lenora. Po několika zaklepáních vletí do místnosti havran, který usedne na bustu Pallady. Havran je symbolem smutku po ztracené lásce. Tím, že se usadí na bustu bohyně moudrosti, naznačuje, že smutek a vzpomínky ovládly básníkovu mysl. Havran opakuje slovo „nevermore“, což znamená, že se láska už nikdy nevrátí. Příběh končí tím, že havran neodletí a na všechny básníkovy otázky odpovídá stále stejně: „Už víckrát ne.“ Konec tedy vyjadřuje rezignaci na život a věčné uvěznění v bolesti ze ztráty.
\end{itemize}

\textbf{Smysl díla:}
\begin{itemize}
\item Dílo ukazuje, jak může člověk po ztrátě blízké osoby propadnout smutku, beznaději a šílenství.
\end{itemize}

\textbf{Pravděpodobný adresát:}
\begin{itemize}
\item Veřejnost
\end{itemize}

\break
\textbf{Kompozice:}
\begin{itemize}
\item Chronologická
\end{itemize}

\textbf{Zařazení do kontextu autorovy tvorby:}
\begin{itemize}
\item Edgar Allan Poe byl považován za melancholického autora a často se ve své tvorbě věnoval tématům ztracené lásky, smrti a smutku. Dílo Havran je jedním z nejznámějších příkladů jeho tvorby.
\end{itemize}

\textbf{Film nebo dramatizace:}
\begin{itemize}
\item Existuje film Havran (2012), ale neviděl jsem ho.
\end{itemize}

\end{tcolorbox}

%------------------------------------------------
\section*{Úryvek k rozboru}

\colorbox{magenta}{\colorbox{teal}{„Proroku,“} dím, „mene tekel, ať jsi pták a nebo z pekel,}\\
synu podsvětí, a přece proroku, pojď hádat mně –\\
statečně, byť opuštěný, žiji zaklet v této zemi,\\
\colorbox{yellow}{dům mám hrůzou obklíčený}, \colorbox{pink}{zda tvá věštba uhádne,}\\
zdali najdu balzám v smrti, \colorbox{pink}{zda tvá věštba uhádne“} –\\
\colorbox{magenta}{\colorbox{lime}{Havran dí}: „Už víckrát ne.“} \\
\\
\colorbox{magenta}{\colorbox{teal}{„Proroku,“} dím, „mene tekel, ať jsi pták a nebo z pekel,}\\
při nebi, jež nad námi je, při Bohu, jenž láká mne,\\
rci té duši, jež žal tají, zdali aspoň jednou v ráji\\
tu, již svatí nazývají Lenora, kdy \colorbox{pink}{přivine},\\
jasnou dívku Lenoru kdy v náručí své \colorbox{pink}{přivine“} –\\
\colorbox{magenta}{\colorbox{lime}{Havran dí}: „Už víckrát ne“}\\

\begin{tcolorbox}[breakable]

\textbf{Atmosféra úryvku:}
\begin{itemize}
\item Napjatá, vážná
\end{itemize}

\textbf{Postavy v úryvku:}
\begin{itemize}
\item Havran
\item Básník
\end{itemize}

\textbf{Charakteristika postav:}

\textbf{Hlavní postava}
\begin{itemize}
\item Havran – není skutečný pták, ale symbol smutku a ztracené lásky Lenory.
\item Básník – hlavní postava, melancholický muž, který prožívá psychické strádání kvůli ztrátě milované ženy.
\end{itemize}

\textbf{Vedlejší postava}
\begin{itemize}
\item -
\item -
\end{itemize}

\textbf{Další postavy díla:}

\begin{itemize}
\item -
\item -
\item -
\end{itemize}

\textbf{Vztahy mezi postavami:}
\begin{itemize}
\item Básník se snaží havrana vyhnat z domova, protože na každou jeho otázku odpovídá „Už víckrát ne“, čímž v něm ještě více prohlubuje beznaděj.
\end{itemize}

\textbf{Zařazení úryvku do děje:}
\begin{itemize}
\item 3/4 díla
\end{itemize}

\textbf{Jazykové prostředky:}
\begin{itemize}
\item \colorbox{magenta}{Refrén}
\item \colorbox{pink}{Epifora}
\item \colorbox{yellow}{Metafora}
\item \colorbox{lime}{Personifikace}
\item \colorbox{teal}{Apostrofa}
\end{itemize}

\end{tcolorbox}

%------------------------------------------------
\section*{Názor na dílo}

Byla to moje první báseň, proto je pro mě její hodnocení trochu těžké. Nejvíce se mi líbila část, kde jsem musel pochopit smysl díla a rozebrat slova, kterým jsem nerozuměl. Celkově pro mě byla báseň Edgara Allana Poea Havran výzvou. Určitě si přečtu i další díla od Poea.

\end{document}
